\begin{abstract}
The long-only mean-variance portfolio, minimising variance for a given return target
subject to linear balance constraints $Bw = c$ (the fully-invested budget constraint
$\mathbf{1}^\top w = 1$ being the leading case) and non-negative weights, is routinely
solved by
assembling an $n\times n$ sample covariance matrix and dispatching it to a
general-purpose quadratic program. Both steps are unnecessary: the covariance matrix
is not a primitive of the problem. It is a computational detour.

The detour is avoidable. The portfolio weights require only a handful of solves
against $\hat\Sigma$ -- in the budget case, $v = \hat\Sigma^{-1}\mathbf{1}$ alone;
conjugate gradients computes them matrix-free, evaluating each product $\hat\Sigma\,v$
as two $O(nT)$ passes over $X$ without ever forming $\hat\Sigma = X^\top X/T$, and the
$p$ balance multipliers are eliminated analytically through a $p \times p$ Schur
complement.
Non-negativity is enforced by a primal-dual active-set outer loop that alternately
removes assets with negative weight (primal step) and re-adds assets that violate
dual feasibility (dual step), restarting the inner solver on the modified active set
each time; the loop and its finite-termination guarantee are developed in full
generality in the companion paper~\cite{schmelzer2026nncg}. The $O(\sqrt\kappa)$ iteration count of CG, where
$\kappa = \kappa(\hat\Sigma)$ is the spectral condition number, contrasts sharply with
the $O(\kappa)$ rate of projected-gradient methods that do not exploit Krylov structure.

Ledoit-Wolf shrinkage~\cite{ledoit2004}, already standard practice for statistical
reasons, also compresses $\kappa$ and reduces solver iterations as a free by-product.
The effect is modest at the statistically optimal intensity
($\alpha_{\mathrm{LW}} \approx 0.017$), because the Frobenius objective is insensitive
to the near-zero eigenvalues that drive ill-conditioning. At the heavier intensities
practitioners apply for out-of-sample stability ($\alpha \approx 0.3$--$0.5$) the
conditioning benefit is decisive and requires no additional preconditioning step.
That practitioners applying heavy shrinkage are simultaneously conditioning their
solver has not been made explicit in the portfolio optimisation literature.
RMT eigenvalue cleaning~\cite{laloux1999,bun2017} provides a more decisive resolution,
decoupling the two objectives; the Woodbury direct solver for the RMT target is
developed in the companion paper~\cite{schmelzer2026rmt}.

We benchmark against general-purpose QP solvers (Clarabel interior-point, OSQP operator-splitting) and a non-Krylov
first-order baseline. On S\&P~500 equity data CG requires 684 iterations without
shrinkage, falling to 82 under heavy LW shrinkage ($\alpha=0.5$). Called head-to-head
against the same solvers through their native APIs the algorithmic gain is roughly
an order of magnitude; against those solvers wrapped in a general modelling layer the
end-to-end gap is larger still, as it also absorbs problem-construction overhead.
CG's runtime grows substantially subquadratically in $n$ while
interior-point methods scale as $O(n^3)$ per iteration, and warm-starting across the
return-tilt sweep traces a 50-point long-only efficient frontier orders of magnitude
faster than a na\"{i}ve loop over a modelling layer.
A rolling-window out-of-sample study on the same universe shows that the heavy-shrinkage
regime which makes the solver fast is also the regime that minimises realised
out-of-sample variance while lowering turnover: the statistical and numerical objectives, far from
conflicting, are reconciled at the intensities practitioners already use.
\end{abstract}
